% !TeX root = ../../Book4.tex
% 纲要标题：### 6.4 实际参数与全局几何
% 纲要位置：几何学计划纲要.md，第 3711 行。
% 正文状态：初稿；尚未正式审核。
% 本次扩充的范围、先修与边界：先修为 IV.5 的双曲统一化与测地几何、IV.6.1 的标记和 Dehn 扭转。主对象为亏格至少二的闭连通定向双曲曲面；选定带标记的裤子分解和缝线，说明长度、扭转的单位、参数数量、有限拼接和恢复。长度趋零只处理选定曲线捏缩的模型，区别 Teichmüller 空间、模空间与稳定退化；不追加完整 Weil–Petersson 几何、Wolpert 公式或一般全局紧化理论。

% 纲要内容开始
% 
% > 补充写作范围与外部输入：先修为 IV.5 的双曲统一化与测地几何、IV.6.1 的标记和 Dehn 扭转。主对象为亏格至少二的闭连通定向双曲曲面；选定带标记的裤子分解和缝线，说明长度、扭转的单位、参数数量、有限拼接和恢复。长度趋零只处理选定曲线捏缩的模型，区别 Teichmüller 空间、模空间与稳定退化；不追加完整 Weil–Petersson 几何、Wolpert 公式或一般全局紧化理论。
% 
% * 局部坐标及自同构的影响
% * 极值长度和 Teichmüller 度量
% * Fenchel–Nielsen 参数的模型
% * 退化、节点和紧化的入口
% 
% * 双曲裤子的边界长度与裤子分解
% * 长度、扭转与标记的选择
% * 沿测地边界的拼接与参数恢复
% * 长度趋零时的标记与退化
% 
% 
% 纲要内容结束

\section{实际参数与全局几何}
\label{g4:sec:0604}

切空间的维数回答了小形变有多少个独立方向；实际参数还须告诉我们如何制造曲面，以及怎样从曲面读回原来的数值。局部复坐标由上一节的切片给出，本节再用双曲裤子的有限拼接构造全局实参数。与此相配，极值长度比较度量形变的大小，长度趋零的模型则说明这些参数怎样走向光滑曲面的边界。

\subsection{局部坐标及自同构的影响}
\label{g4:subsec:0604:01}

固定闭连通定向曲面 \(S\)，亏格 \(g\ge2\)。\cref{g4:thm:0603:local-parameters}把标记黎曼曲面 \([X,f]\) 的一个邻域表示为 \(\mathbb C^{3g-3}\) 中的开集：有限维参数来自 Beltrami 类，光滑坐标变化的方向由 \(\bar\partial\) 像消去。这里的参数点仍带标记；忘记标记以后，还须识别由曲面自同构及映射类作用产生的等价。

\begin{proposition}
	\label{g4:prop:0604:finite-automorphisms}
	设 \(X\) 为亏格 \(g\ge2\) 的紧连通黎曼曲面，则 \(\operatorname{Aut}(X)\) 有限。
\end{proposition}
\begin{proof}
	每个自同构保持 \(X\) 的 Poincaré 度量，这是\cref{g4:thm:0504:poincare-complete}对映射及其逆同时应用的结果。于是自同构族在紧度量空间上等度连续；取有限网格并作对角选择，任意序列都存在一致收敛子列。对逆映射作同样选择，极限互为逆，因而是等距同胚。局部取相对紧全纯坐标片，一致收敛与 Cauchy 公式又说明极限及其逆全纯。因此自同构族在一致拓扑下紧。

	若它无限，可以由上述紧性构造互不相同的 \(F_n\to\operatorname{id}\)。记 \(\varepsilon_n=\max_{p\in X}d_X(F_n(p),p)>0\)。在有限张较小坐标片内，\((F_n(z)-z)/\varepsilon_n\) 局部一致有界；由 Montel 定理同时取子列，使它们局部一致收敛。若换坐标 \(w=h(z)\)，Taylor 展开给出
	\[
		\frac{h(F_n(z))-h(z)}{\varepsilon_n}
		=h'(z)\frac{F_n(z)-z}{\varepsilon_n}+O(\varepsilon_n).
	\]
	故极限按向量场的变换律粘合为一个全纯向量场 \(V\)。选取实现最大位移的点并再取收敛子列，局部度量与坐标距离的双边比较说明 \(V\ne0\)。但 \(\deg T_X=2-2g<0\)，由\cref{g4:thm:0305:high-degree-sections}，负次线丛没有非零全纯截面，矛盾。因此自同构群有限。
\end{proof}

局部模空间确实只需作这个有限商。为说明不会有其他映射类从很远处回到同一小邻域，取局部参数 \(a\) 对应的光滑双曲度量 \(h_a\)，它在固定标记曲面上随 \(a\) 光滑变化。假如有 \(a_n,b_n\to0\) 及等价 \(F_n:(S,h_{a_n})\to(S,h_{b_n})\)，则这些等距映射及其逆在固定背景度量下具有统一的 Lipschitz 界。上面的紧性论证给出子列收敛于 \((S,h_0)\) 的自同构 \(F\)。还可把收敛提高到光滑收敛：在一点选取微分的收敛子列，局部等距满足 \(F_n(\exp_p^{h_{a_n}}v)=\exp_{F_n(p)}^{h_{b_n}}(\mathrm{d}F_n|_p v)\)，而测地线方程随初值及光滑度量光滑变化。有限个小片逐次延续便得到全曲面上的 \(C^1\) 收敛。当 \(F^{-1}F_n\) 在 \(C^1\) 意义下足够接近恒等时，以短测地线写成 \(p\mapsto\exp_pV_n(p)\)，并用 \(p\mapsto\exp_p(tV_n(p))\)、\(0\le t\le1\) 插值。每一步的微分都可逆，紧性使它为有限覆盖，同伦度为 \(1\) 又使它为微分同胚，故这确实是一条同位。因此 \(F_n\) 最终属于 \(F\) 的映射类。

\ProseParagraphBreak

由于 \(\operatorname{Aut}(X)\) 有限，可缩小局部参数邻域，并对有限群取共同不变的邻域，使忘记标记的局部空间恰是该邻域模有限稳定子。作用是全纯的，因为标记变换与局部切片的坐标转移均为全纯。有限群作用还可局部线性化：若各 \(A\) 固定原点，置
\[
	\Psi(z)=\frac1{|G|}\sum_{A\in G}(\mathrm{d}A_0)^{-1}A(z).
\]
则 \(\mathrm{d}\Psi_0=I\)，且 \(\Psi(Az)=\mathrm{d}A_0\Psi(z)\)，逆函数定理给出所需坐标。因此局部模空间可写成一个有限线性群商。该表示可能有核；非平凡稳定子也不自动造成商空间的奇点。

\subsection{极值长度和 Teichmüller 度量}
\label{g4:subsec:0604:02}

设 \(F:X\to Y\) 为保向 \(C^1\) 微分同胚。在任一点取共形正交坐标，令其微分的两个奇异值为 \(s_1\ge s_2>0\)，定义最大伸缩比
\[
	K(F)=\sup_X\frac{s_1}{s_2}\ge1.
\]
这个量不受源、目标共形密度的改变影响；在复坐标中它等于 \((1+\|\mu_F\|_\infty)/(1-\|\mu_F\|_\infty)\)。当 \(K(F)=1\) 时，微分是复线性的，故 \(F\) 双全纯。

\begin{proposition}
	\label{g4:prop:0604:extremal-qc}
	设 \(F:X\to Y\) 是保向 \(C^1\) 微分同胚，且 \(K(F)\le K<\infty\)。对 \(X\) 上任一曲线族 \(\Gamma\)，有
	\[
		K^{-1}\operatorname{EL}_X(\Gamma)
		\le\operatorname{EL}_Y(F\Gamma)
		\le K\operatorname{EL}_X(\Gamma).
	\]
\end{proposition}
\begin{proof}
	给定目标可容许密度 \(\sigma\)，在源上取 \(\rho=(\sigma\circ F)s_1\)。链式法则说明每条曲线都有
	\[
		L_\sigma(F\gamma)\le L_\rho(\gamma).
	\]
	另一方面，\(s_1^2\le K s_1s_2=KJ_F\)，面积换元给出
	\[
		A_\rho\le K\int_X(\sigma\circ F)^2J_F\,\mathrm{d}A_X
		=KA_\sigma.
	\]
	因此 \(L_\sigma(F\Gamma)^2/A_\sigma\le K\operatorname{EL}_X(\Gamma)\)。对 \(\sigma\) 取上确界得到右边不等式；对 \(F^{-1}\) 应用同样论证得到左边。背景共形度量的选择不影响计算。
\end{proof}

一般准共形映射只有弱导数，沿某条指定曲线未必绝对连续；因而不能逐条直接套用上述证明。本节对这种映射使用环域版本，证明中只沿几乎每条平行圆周积分。这里所需的逆映射、复合、零测集保持及面积换元性质，采用准共形解析理论的标准结论；下面的证明也是 Glutsyuk 作者稿引理3.7的圆柱形式。\cite[作者稿第12页，引理3.7；群性质与零测集保持见第3.4节]{Glutsyuk2008Uniformization}

\begin{proposition}
	\label{g4:prop:0604:annulus-qc}
	设 \(A_1,A_2\) 为有限模量的共形圆环域，\(F:A_1\to A_2\) 为 \(K\)-准共形同胚，则
	\[
		K^{-1}\operatorname{Mod}(A_1)
		\le\operatorname{Mod}(A_2)
		\le K\operatorname{Mod}(A_1).
	\]
	因此环域的连接族与核心族也分别满足同一极值长度比较。
\end{proposition}
\begin{proof}
	利用共形坐标，将 \(A_j\) 写成平坦圆柱
	\[
		(0,L_j)\times(\mathbb R/2\pi\mathbb Z),
		\qquad \operatorname{Mod}(A_j)=L_j/(2\pi).
	\]
	记源坐标为 \((x,\theta)\)。局部 \(W^{1,2}\) 性说明对几乎所有 \(x\)，\(F(x,\theta)\) 关于 \(\theta\) 绝对连续。这个闭路在目标圆柱的绕行数为 \(\pm1\)，所以其平坦长度至少为 \(2\pi\)。几乎处处的伸缩比界给出
	\[
		|\partial_\theta F|^2\le KJ_F.
	\]
	先在 \(\varepsilon<x<L_1-\varepsilon\) 积分，再令 \(\varepsilon\downarrow0\)，由面积换元、Fubini 定理\footnote{Guido Fubini（1879--1943），意大利数学家，研究分析、微分几何与重积分。}与 Cauchy--Schwarz 不等式，
	\[
		2\pi L_2=\int_{A_1}J_F
		\ge\frac1K\int_0^{L_1}\int_0^{2\pi}|\partial_\theta F|^2
		\,\mathrm{d}\theta\,\mathrm{d}x
		\ge\frac{2\pi L_1}{K}.
	\]
	这给出左边不等式。逆映射仍为 \(K\)-准共形，交换源、目标便得右边。最后由\cref{g4:thm:0505:annulus-modulus}，两族的极值长度分别等于模量及其倒数，所以得到所述比较。
\end{proof}

标记规定了允许比较的映射同伦类。\term[Teichmullerjuli]{Teichmüller 距离}[Teichmüller distance]定义为
\begin{equation}
	\label{g4:eq:0604:teich-distance}
	d_{\mathrm T}([X,f],[Y,h])
	=\frac12\inf_{F\simeq h\circ f^{-1}}\log K(F),
\end{equation}
其中下确界取遍该类的准共形同胚；同伦类与光滑同位标记的衔接由\cref{g4:prop:0601:surface-isotopy-interface}保证。\(K(F)\) 是几乎处处伸缩比的本质上确界。取一个光滑标记比较映射，紧性便给出有限的 \(K(F)\)，所以距离值有限。逆映射保持 \(K\)，复合满足 \(K(GF)\le K(G)K(F)\)，因而对称性与三角不等式直接来自伸缩比的运算。

\ProseParagraphBreak

严格正定性及完备性涉及准共形映射族的全局紧性。这里把以下结论列为外部输入：式\eqref{g4:eq:0604:teich-distance}确实是 \(\mathcal T(S)\) 上的完备距离，并由其通常复流形切空间上的连续范数通过曲线长度产生，因此诱导前一节的局部参数拓扑。Hamenstädt 的综述在第3节陈述了这一结论；此处引用的是作者公开稿的该节，页码与正式文集版本不同。\cite[作者公开稿第3节，第234--235页]{Hamenstadt2006Mapping}
\cref{g4:thm:0602:qc-compactness}给出了这类论证中的基本解析输入：固定伸缩界的规范化映射不能任意坍缩，伸缩比趋于 \(1\) 时其极限恢复共形性。本节后面的长度--扭转参数化将由有限双曲拼接另行证明。

\begin{example}
	\label{g4:exm:0604:annulus-distortion}
	若两个有限圆环域的模量为 \(M_1,M_2\)，则任一准共形同胚的伸缩比至少为
	\[
		\max\{M_2/M_1,M_1/M_2\}.
	\]
	在上述圆柱坐标中，映射
	\((x,\theta)\mapsto(L_2x/L_1,\theta)\)
	达到这个下界，因为它两个方向的伸缩分别为 \(L_2/L_1\) 与 \(1\)。因此环域之间按式\eqref{g4:eq:0604:teich-distance}取下确界时，数值为
	\(\tfrac12|\log(M_2/M_1)|\)。
	这是一个完全可计算的参数模型：很长的圆柱与较短圆柱之间必有明显的非共形伸缩。
\end{example}

\subsection{Fenchel–Nielsen 参数的模型}
\label{g4:subsec:0604:03}

在 \(S\) 上固定一组互不相交、互不同伦且本质的简单闭曲线 \(c_1,\ldots,c_N\)，把 \(S\) 切成三边界球面；每个这样的分块称为一个\term[shuangqukuzi]{双曲裤子}[hyperbolic pair of pants]的拓扑模型。为各块标号，并在每块上选定连接不同边界的三条缝线的同位类。带这些辅助资料的分解将在下面得到几何实现。

\ProseParagraphBreak

每个曲线 \(c_i\) 将给出两个实数：其双曲测地代表的长度 \(\ell_i>0\)，以及把两侧裤子拼回时沿该测地线平移的有向距离 \(\tau_i\in\mathbb R\)。后者的单位是长度；一次完整 Dehn 扭转改变 \(\tau_i\) 一个 \(\ell_i\)。若改用圈数 \(u_i=\tau_i/\ell_i\)，一次完整扭转就改变 \(u_i\) 一个整数。这就是\term[FenchelNielsencanshu]{Fenchel--Nielsen 参数}[Fenchel--Nielsen coordinates]\footnote{Werner Fenchel（1905--1988），德国出生、后在丹麦工作的数学家，研究凸几何与非 Euclid 几何。}的模型：
\[
	(\ell_1,\tau_1,\ldots,\ell_N,\tau_N)
	\in(\mathbb R_{>0}\times\mathbb R)^N.
\]
具体的六边形构造、扭转的提升选择及恢复将在本节后四个小节证明。

\begin{proposition}
	\label{g4:prop:0604:pants-count}
	闭连通定向亏格 \(g\ge2\) 的曲面存在裤子分解；每个裤子分解含 \(2g-2\) 条裤子和 \(N=3g-3\) 条切割曲线。因此长度与扭转总共提供 \(6g-6\) 个实参数。
\end{proposition}
\begin{proof}
	由第一卷命题6.4.5，可以在标准手柄模型中工作。先沿每个手柄的一条非分离曲线切开，得到有 \(2g\) 个边界的球面。每次用一条圈从这个球面分出两个边界与一个新切口组成的裤子，直到余块也是裤子，就得到所需分解。这些切线均本质且不平行；把最初的成对边界重新粘合，便回到 \(S\)。

	由第一卷命题6.5.2，沿边界圆周粘合时 Euler 示性数相加。对任一裤子分解，设裤子数为 \(P\)。每条裤子的 Euler 示性数为 \(-1\)，沿圆周粘合不改变其和，所以 \(-P=2-2g\)，即 \(P=2g-2\)。每条裤子有三条边界，而每条切线贡献两侧，故 \(3P=2N\)，于是 \(N=3g-3\)。同一裤子的两条边界互相拼合时仍计两侧，这个计数同样成立。
\end{proof}

\subsection{退化、节点和紧化的入口}
\label{g4:subsec:0604:04}

参数 \(\ell_i>0\) 描述光滑曲面；令 \(\ell_i\to0\) 会离开这个参数域。在共形坐标中，模型是把两端按 \(zw=t\)、\(t\ne0\) 接成环带，\(t\to0\) 时得到两条分支 \(zw=0\)。第五章已计算其模量
\[
	\operatorname{Mod}=\frac1{2\pi}\log\frac{c^2}{|t|}\longrightarrow\infty
\]
及两侧固定截口观察下的尖点极限。相应的结点不是某个光滑黎曼曲面的普通坐标点，而是两条局部分支相遇的位置。

\ProseParagraphBreak

标记与忘记标记会给出不同的边界行为。固定长度后反复作 Dehn 扭转，可能只改变标记而不改变底层曲面；缩短被扭转的圈则改变了共形几何。这两种操作不能用同一个“参数发散”代替。\secref{g4:sec:0605}将在有序四点球面上通过交比把三个碰撞方向写成具体的结点族；本节则在最后一个小节固定被捏缩的裤子曲线，追踪长度、缝线与扭转。一般紧化还要说明允许哪些奇异曲线、如何比较全部退化族以及标记的遗忘方式。

\subsection{双曲裤子的边界长度与裤子分解}
\label{g4:subsec:0604:05}

双曲裤子的三条边界都是闭测地线。把裤子沿三条最短缝线切开，就应得到两个全等的直角六边形；边界长的一半成为交替的三条边。下面先直接构造这种六边形，借此同时证明任意三个正边界长度都能实现，而且没有隐藏的额外参数。

\begin{lemma}
	\label{g4:lem:0604:right-hexagon}
	给定 \(a,b,c>0\)，存在一个紧凸双曲直角六边形，使三条交替边长为 \(a,b,c\)，并且它在允许反向的等距下唯一。连接 \(a,b\) 两边的边长 \(x\) 满足
	\begin{equation}
		\label{g4:eq:0604:hexagon-cosine}
		\cosh x=\frac{\cosh c+\cosh a\cosh b}{\sinh a\sinh b}.
	\end{equation}
	其余两条边由循环置换得到。
\end{lemma}
\begin{proof}
	采用第一卷定理10.1.3的双曲面距离与测地线公式。在 Lorentz 模型\footnote{Hendrik Antoon Lorentz（1853--1928），荷兰理论物理学家，研究电磁理论、电子理论与运动物体的变换。} \(\langle u,v\rangle=-u_0v_0+u_1v_1+u_2v_2\)、\(\mathbb H^2=\{\langle u,u\rangle=-1,u_0>0\}\) 中，取两条有公垂线的测地线，垂足为
	\[
		p_1=(1,0,0),\qquad p_2=(\cosh x,\sinh x,0),
	\]
	并沿二者相同一侧的单位切向量 \((0,0,1)\) 分别走距离 \(a,b\)，到达
	\[
		q_1=(\cosh a,0,\sinh a),\quad
		q_2=(\cosh b\cosh x,\cosh b\sinh x,\sinh b).
	\]
	在 \(q_1,q_2\) 处作垂直于原两条测地线的新测地线 \(h_1,h_2\)。它们的单位法向量分别为
	\[
		n_1=(\sinh a,0,\cosh a),\quad
		n_2=(\sinh b\cosh x,\sinh b\sinh x,\cosh b).
	\]
	因而 \(\langle n_1,n_2\rangle=\cosh a\cosh b-\sinh a\sinh b\cosh x\)。式\eqref{g4:eq:0604:hexagon-cosine}的右边大于 \(1\)，因为分子减分母为 \(\cosh c+\cosh(a-b)>0\)，所以它唯一确定 \(x>0\)。按此选择有 \(\langle n_1,n_2\rangle=-\cosh c<-1\)。两个单位空间向量张成的平面因此为 Lorentz 平面，正交补为正定直线；在它们张成的平面中取单位双曲向量，直接归约到 \((0,1,0)\) 和 \((\sinh c,-\cosh c,0)\) 的模型，便知两测地线不相交，其唯一公垂线长为 \(c\)。

	还须保证公垂足在所需的正向射线上。沿 \(h_1\) 向第二条原测地线走，参数式为
	\[
		r_1(t)=\cosh t\,q_1+\sinh t\,(0,1,0).
	\]
	令 \(A=-\langle q_1,n_2\rangle=(\cosh a\cosh c+\cosh b)/\sinh a\)，\(B=\sinh b\sinh x\)。则 \(A,B>0\)，且代入上式可算得 \(A^2-B^2=\sinh^2c\)。距离 \(h_2\) 的垂足由 \(\tanh t=B/A\) 确定，所以 \(t>0\)，并有
	\[
		\cosh t=\frac{\cosh b+\cosh a\cosh c}{\sinh a\sinh c}.
	\]
	另一侧完全对称，得到另一个正边长。两足还位于原两测地线之间：例如上面的 \(\tanh t<\cosh a\tanh x\) 正是 \(r_1(t)\) 尚未越过第二条原测地线的条件，代入化为 \(\cosh c<\cosh^2a\cosh c\)。两原测地线围成的带及各垂线的内半平面均测地凸，故这六条边按次序围出一个紧凸六边形，六角都是直角。

	反过来，任一直角六边形可由等距把前述两条边及其间公垂边放到上述位置。相邻垂线的法向内积必为 \(-\cosh c\)，从而强迫式\eqref{g4:eq:0604:hexagon-cosine}；\(x\) 及所有垂线随之确定，其余两条边也唯一。这证明存在与唯一，且给出了全部边长的计算。
\end{proof}

\begin{theorem}
	\label{g4:thm:0604:pants-lengths}
	给定三个正数 \(\ell_1,\ell_2,\ell_3\)，存在边界依次标号的紧双曲裤子，其测地边界长为这些数；它在保向且保持边界标号的等距下唯一。连接第 \(i,j\) 条边界的最短缝线长 \(s_{ij}\) 满足
	\[
		\cosh s_{ij}
		=\frac{\cosh(\ell_k/2)+\cosh(\ell_i/2)\cosh(\ell_j/2)}
		{\sinh(\ell_i/2)\sinh(\ell_j/2)},\qquad \{i,j,k\}=\{1,2,3\}.
	\]
\end{theorem}
\begin{proof}
	将两个交替边为 \(\ell_1/2,\ell_2/2,\ell_3/2\) 的六边形沿另外三边反射粘合。测地边的反射是局部双曲等距，内点两侧的度量因此光滑；在各端点，两直角合成平角，所以留下的边界也是光滑测地线。拓扑上得到三边界球面，且三条边界长正好为 \(\ell_1,\ell_2,\ell_3\)。

	反向考虑任一双曲裤子。不同边界之间的距离在紧性下达到；最短弧的内部是测地线，端点自由变分使它与两边界垂直。它不能在途中接触第三条测地边界，否则相切后由测地线唯一性不能再离开。三对边界的最短弧彼此不交：若两弧从同一个边界到达一个交点，交换从该边界到交点的两段，较短的一段会缩短另一条弧；若两段等长，交换后在交点形成的角又能严格缩短。重叠一段后再分开也违反测地线初值唯一性。

	三条不交弧把裤子切成两个圆盘，每个圆盘的边界有六条边且六角为直角。这个拓扑事实可在球面去三个开盘的模型中沿弧逐条切开核对。两六边形的三条交替缝线边分别相等，故由上一引理唯一性，两六边形全等。于是各边界被分成相等两半，边长恰为 \(\ell_i/2\)。引理便同时给出唯一性与缝线公式；反射对称也说明上述三条缝线由几何唯一确定。
\end{proof}

% !TeX root = ../Book4.tex
\begin{figure}[htbp]
	\centering
	\begin{tikzpicture}[font=\small,scale=1.05]
		\fill[GeometryBlueLight] (0,0) ellipse (3 and 1.65);
		\fill[white] (-1.2,0) circle (0.62);
		\fill[white] (1.2,0) circle (0.62);
		\draw[GeometryGray,thick] (0,0) ellipse (3 and 1.65);
		\draw[GeometryGray,thick] (-1.2,0) circle (0.62);
		\draw[GeometryGray,thick] (1.2,0) circle (0.62);
		\draw[GeometryBlue,very thick] (-3,0)--(-1.82,0);
		\draw[GeometryBlue,very thick] (-0.58,0)--(0.58,0);
		\draw[GeometryBlue,very thick] (1.82,0)--(3,0);
		\node[above] at (-2.4,0.05) {\(s_{13}\)};
		\node[above] at (0,0.05) {\(s_{12}\)};
		\node[above] at (2.4,0.05) {\(s_{23}\)};
		\node at (-1.2,0) {\(C_1\)};
		\node at (1.2,0) {\(C_2\)};
		\node[above] at (0,1.65) {\(C_3\)};
		\node at (0,1) {上侧六边形};
		\node at (0,-1) {下侧六边形};
		\node[below] at (-1.2,-0.62) {\(\ell_1/2\)};
		\node[below] at (1.2,-0.62) {\(\ell_2/2\)};
		\node[below] at (0,-1.65) {\(\ell_3/2\)};
	\end{tikzpicture}
	\caption[裤子的三条缝线与六边形分解]{三条蓝色缝线分别连接不同边界，把裤子分成上下两个直角六边形。每个六边形的三条边界弧长为 \(\ell_1/2,\ell_2/2,\ell_3/2\)，另三边为缝线。图为平面拓扑示意，所画圆周和长度比例不表示双曲测量。}
	\label{g4:fig:0604:pants-seams}
\end{figure}


\cref{g4:fig:0604:pants-seams}把三条缝线与两个六边形的对应画出；图中的圆周只表示边界的拓扑位置，六角的直角性来自双曲度量。

对闭双曲曲面，所选裤子曲线都能同时换成互不相交的测地线。其原因可以从万有覆盖看清：一个非平凡自由同伦类对应一个覆盖变换。紧商可由有限个均匀覆盖的小圆盘覆盖，取这些盘的 Lebesgue 数便给出非平凡覆盖变换位移的正下界。由第一卷定理10.2.1，该变换不能是抛物型；自由作用又排除椭圆型，所以它有唯一双曲轴，投影得到该类的闭测地线。原简单曲线的一条提升由一个紧基本弧在该覆盖变换下的像拼成，因而与双曲轴保持有界距离；它是连接轴的两个理想端点的嵌入直线。不同提升不相交。若相应轴相交，其端点在理想圆周上交错，第一卷定理6.2.3的平面分离迫使原提升也相交，矛盾。因此测地代表仍简单；对两条不同且不交的曲线作同样比较，测地代表也不交。由\cref{g4:prop:0601:surface-isotopy-interface}，简单闭曲线的自由同伦与同位相容，所以这些测地线代表原来指定的切割同位类。切开以后正是上一定理描述的双曲裤子。

\subsection{长度、扭转与标记的选择}
\label{g4:subsec:0604:06}

边界长为 \(\ell\) 的测地圆周，可用弧长坐标 \(t\in\mathbb R/\ell\mathbb Z\) 表示。给每条裤子边界指定其中一条缝线的垂足作为零点，并固定与参考曲线方向相容的坐标方向。两侧边界要反向粘合以保持曲面定向；若各用本侧诱导的边界方向，拼接式写为
\[
	t_-=-t_++\tau\pmod\ell.
\]
若两侧改用同一个参考曲线方向，则它就是一个平移。几何上的拼接只看见 \(\tau\bmod\ell\)，而标记还记录绕行的整数次数。

\ProseParagraphBreak

为把这个整数固定下来，在参考曲面的每个切割环带中选一条横穿弧，连接两侧指定的缝线。将标记送来的这条弧与几何缝线对齐端点，在环带的条带覆盖中提升；规定左端弧长坐标为零后，右端坐标的实数值就是 \(\tau\)。更换弧的相对同伦类一个正整圈，会使 \(\tau\) 增加 \(\ell\)。所以 \(\tau\in\mathbb R\) 由曲面、标记及固定的缝线与横穿弧共同确定，不能在保留标记时把它任意约到 \([0,\ell)\)。

\ProseParagraphBreak

Dehn 扭转的符号可按这一条带坐标固定：称使提升右端增加 \(\ell\) 的扭转为本参数的正扭转；其对标记的作用为 \((\ell,\tau)\mapsto(\ell,\tau+\ell)\)。若映射类作用使用预合成逆映射的约定，同一个群元素在参数上的符号也相应反转。零点换用另一条缝线通常增加半个边界长度，改变横穿弧则增加整倍边界长度；这些都是辅助选择的改变，不是曲面新添了连续参数。

\subsection{沿测地边界的拼接与参数恢复}
\label{g4:subsec:0604:07}

\begin{theorem}
	\label{g4:thm:0604:fenchel-nielsen}
	固定亏格 \(g\ge2\) 的闭连通定向参考曲面 \(S\)，以及上述带标号的裤子分解、缝线和横穿弧，令 \(N=3g-3\)。长度与实扭转给出一一对应
	\[
		\mathcal T(S)\longrightarrow
		(\mathbb R_{>0}\times\mathbb R)^N,
		\qquad [X,f]\longmapsto(\ell_i,\tau_i)_{i=1}^N.
	\]
	在局部 Teichmüller 拓扑与右边的通常拓扑下，它及其逆均连续。
\end{theorem}
\begin{proof}
	先从任意一组参数构造曲面。每条裤子的三条边界长由相邻 \(\ell_i\) 指定，\cref{g4:thm:0604:pants-lengths}给出它的唯一模型。沿所有配对边界按 \(t_-=-t_++\tau_i\) 粘合。每条测地边界附近有法向坐标，度量为
	\[
		\mathrm{d}r^2+\cosh^2r\,\mathrm{d}t^2,
	\]
	其中边界为 \(r=0\)。该式可在上半平面轴旁直接计算，正是\cref{g4:eq:0504:cylinder-metric}的局部版本。反向配对以后以有符号的 \(r\) 合并两侧，\(\cosh r\) 为偶函数，度量光滑通过接缝，且曲率仍为 \(-1\)。有限块粘合成紧定向曲面；其全纯结构由局部双曲图给出。

	标记在每条裤子上由缝线和边界标号确定，再在各环带中用所指定实位移的条带提升接起来。相差整周的位移给同一个未标记的边界配对，却使横穿弧多绕一圈，因而按定义改变标记的 Dehn 扭转数。这说明全部实数 \(\tau_i\) 都有实现。

	反向给定标记曲面，以统一化的 Poincaré 度量将各 \(c_i\) 化为测地代表并切开。边界长给出 \(\ell_i\)，每条裤子的唯一模型确定其缝线位置；标记的横穿弧在条带覆盖中的终点给出 \(\tau_i\)。因此构造后再读取参数会恢复原数据。若两张标记曲面参数相同，各裤子的保向等距在缝线、边界及横穿弧上相容，遂拼成全局等距，其标记差固定所选切割弧；在剩下的圆盘中以及扭转数为零的环带中，\cref{g4:lem:0601:relative-disk-annulus}将它相对同位为恒等。它们代表同一个 Teichmüller 点。等距是双全纯的；反过来双全纯等价保持 Poincaré 度量与标记数据，所以对应在定义的等价类上成立。

	最后核对连续性。六边形的坐标、边长和垂足公式在全部 \(\ell_i>0\) 的区域光滑。用缝线、法向领圈及内部有限三角片识别邻近裤子，得到随参数连续且局部光滑的标记度量；领圈中的扭转由一个固定截断函数乘 \(\tau_i/\ell_i\) 实现。因此逆构造在局部切片拓扑中连续。反向在局部切片中选取连续依赖参数的光滑双曲度量。各固定类的测地代表长度局部有界；按弧长参数化后，由紧性与测地线方程，其任一序列都有光滑收敛子列。极限仍在原自由同伦类，并为测地线，唯一性使它就是极限度量的代表；因此全序列收敛。裤子内垂直缝线同样由最短性、紧性和唯一性连续变化。由此长度和扭转模 \(\ell_i\) 连续；横穿弧固定的整数提升在足够小邻域中不跳变，所以实扭转也连续。
\end{proof}

这个参数化的实维数为 \(2(3g-3)\)，与\cref{g4:cor:0603:dimension}的复维数 \(3g-3\) 相符。长度和扭转是实几何坐标，不能逐对把 \(\ell_i+\mathrm{i}\tau_i\) 直接当作前一节的全纯坐标；两种坐标分别服务于不同的计算。

\subsection{长度趋零时的标记与退化}
\label{g4:subsec:0604:08}

固定裤子分解中的一条曲线 \(c\)，让它的长度 \(\ell\downarrow0\)，并把其余切割曲线的长度限制在正的紧区间中。缝线公式立刻显示通向这条边界的缝线越来越长。例如，若邻近裤子的另外两条边界长 \(b,c_0\) 保持正值，则
\[
	\cosh s
	=\frac{\cosh(c_0/2)+\cosh(\ell/2)\cosh(b/2)}
	{\sinh(\ell/2)\sinh(b/2)},
\]
所以 \(s=\log(1/\ell)+O(1)\)。若切线两侧来自同一条裤子，便有两条边界同时长为 \(\ell\)：连接这两条边界的缝线长为 \(2\log(1/\ell)+O(1)\)，它们到第三条正长度边界的缝线仍为 \(\log(1/\ell)+O(1)\)，直接代入同一公式即可得到。原来有限长度的边界因而远离固定的内部部分；将未收缩部分固定在模型中，相关六边形顶点趋于理想边界，得到以一个或两个尖点代替相应测地边界的裤子。

\ProseParagraphBreak

对指定的嵌入双曲环带，这个变化有更精确的尺度。\cref{g4:prop:0505:truncated-cylinder}计算了核心长 \(\ell\)、两侧截口长固定为 \(a>\ell\) 的对称模型：其模量为
\[
	M_\ell=\frac2\ell\arccos(\ell/a)
	=\frac\pi\ell-\frac2a+O(\ell^2).
\]
从任一截口向内量距离 \(s\)，度量系数为
\[
	\mathrm{d}s^2+
	\bigl(a\cosh s-\sqrt{a^2-\ell^2}\sinh s\bigr)^2\,\mathrm{d}x^2,
	\qquad x\in\mathbb R/\mathbb Z.
\]
它在每个固定宽度区域上光滑趋于 \(\mathrm{d}s^2+a^2e^{-2s}\mathrm{d}x^2\)。这精确表达了两侧看到两个尖点，而中间核心退到无限距离处的现象。对当前裤子拼接族，可以在理想六边形极限的尖点内先选足够短的截口，再取邻近参数的对应截口，得到这样的嵌入端部；这里处理的是指定拼接族，未对所有短闭测地线宣称同一个全局领圈宽度。

\ProseParagraphBreak

标记仍记住这条被捏缩的曲线及横穿它的弧。实扭转 \(\tau\) 与圈数 \(u=\tau/\ell\) 的行为须分别记录：即使 \(\tau\) 有界，\(u\) 也可能无界；而作一次 Dehn 扭转改变 \(\tau\) 的量本身趋于零。忘记标记以后，\(u\) 的整数部分被相应映射类识别。将两个尖点补成结点时，原来的边界圆不再存在，沿该圆的扭转也不再作为结点曲线本身的一个参数保留；邻近光滑族的粘合参数 \(t\) 仍通过其辐角记录相对旋转。

\ProseParagraphBreak

若同时捏缩若干条互不相交的裤子曲线，可在各配对边界分别进行上述构造。切开后每个分块具有负 Euler 示性数；把尖点按原配对补为结点，所得分支满足稳定性条件：球面分支至少有三个特殊点，环面分支至少有一个。其余正长度和扭转参数描述未捏缩部分。这个有限模型解释了光滑参数域怎样邻接结点曲线，并为下一节四点球面的三个具体边界提供度量图像。
% 纲要条目排为编号小节标题，并进入目录与书签；说明段仍保留于纲要注释中。
